%%%%%%%%%%%%%%%%%%%%%%%%%%%%%%%%%%%%%%%%%
% Journal Article
% LaTeX Template
% Version 1.3 (9/9/13)
%
% This template has been downloaded from:
% http://www.LaTeXTemplates.com
%
% Original author:
% Frits Wenneker (http://www.howtotex.com)
%
% License:
% CC BY-NC-SA 3.0 (http://creativecommons.org/licenses/by-nc-sa/3.0/)
%
%%%%%%%%%%%%%%%%%%%%%%%%%%%%%%%%%%%%%%%%%

%----------------------------------------------------------------------------------------
%	PACKAGES AND OTHER DOCUMENT CONFIGURATIONS
%----------------------------------------------------------------------------------------

\documentclass[twoside]{article}

%\documentclass{aastex}  % version 5.0 or prior
%\usepackage{natbib}



\usepackage{graphicx}
\usepackage{lipsum} % Package to generate dummy text throughout this template
%\usepackage[sc]{mathpazo} % Use the Palatino font
\usepackage[T1]{fontenc} % Use 8-bit encoding that has 256 glyphs
\linespread{1.05} % Line spacing - Palatino needs more space between lines
\usepackage{microtype} % Slightly tweak font spacing for aesthetics

\usepackage[margin=1in,columnsep=20pt]{geometry} % Document margins
\usepackage{multicol} % Used for the two-column layout of the document
\usepackage[hang, small,labelfont=bf,up,textfont=it,up]{caption} % Custom captions under/above floats in tables or figures
\usepackage{booktabs} % Horizontal rules in tables
\usepackage{float} % Required for tables and figures in the multi-column environment - they need to be placed in specific locations with the [H] (e.g. \begin{table}[H])
\usepackage{hyperref} % For hyperlinks in the PDF
\usepackage{subcaption}

\usepackage{lettrine} % The lettrine is the first enlarged letter at the beginning of the text
\usepackage{paralist} % Used for the compactitem environment which makes bullet points with less space between them
\usepackage{amsmath}
\usepackage{abstract} % Allows abstract customization
\renewcommand{\abstractnamefont}{\normalfont\bfseries} % Set the "Abstract" text to bold
\renewcommand{\abstracttextfont}{\normalfont\small\itshape} % Set the abstract itself to small italic text

\usepackage{titlesec} % Allows customization of titles
%\renewcommand\thesection{\Roman{section}} % Roman numerals for the sections
%\renewcommand\thesubsection{\Roman{subsection}} % Roman numerals for subsections
%\renewcommand\thesubsubsection{\Alph{subsubsection}} % Roman numerals for subsections
\titleformat{\section}[block]{\Large\scshape}{\thesection}{1em}{} % Change the look of the section titles
\titleformat{\subsection}[block]{\large}{\thesubsection}{1em}{} % Change the look of the section titles
\titleformat{\subsubsection}[block]{}{\thesubsubsection}{1em}{} % Change the look of the section titles

\usepackage{fancyhdr} % Headers and footers
\pagestyle{fancy} % All pages have headers and footers
\fancyhead{} % Blank out the default header
\fancyfoot{} % Blank out the default footer
\fancyhead[C]{Montana State University \quad $\bullet$ \quad CSCI 466 Artificial Intelligence \quad $\bullet$ \quad Group 21} % Custom header text
\fancyfoot[RO,LE]{\thepage} % Custom footer text

\newcommand{\ve}[1]{\boldsymbol{\mathbf{#1}}}

%----------------------------------------------------------------------------------------
%	TITLE SECTION
%----------------------------------------------------------------------------------------

\title{\vspace{-15mm}\fontsize{24pt}{10pt}\selectfont\textbf{CSCI 446 Artificial Intelligence \\ Project 1 Design Report} \\[-2mm]} % Article title
\date{\today}
\author{
\large
\textsc{Roy Smart} \and \textsc{Nevin Leh} \and \textsc{Brian Marsh}\\[2mm] % Your name
}


%----------------------------------------------------------------------------------------

\begin{document}

\maketitle % Insert title

\thispagestyle{fancy} % All pages have headers and footers

%\begin{abstract}
%We present a novel way of performing MOSES data inversions using a
%\end{abstract}

%----------------------------------------------------------------------------------------
%	ARTICLE CONTENTS
%----------------------------------------------------------------------------------------

\begin{multicols}{2} % Two-column layout throughout the main article text

\section{Introduction}
The \textit{Graph Coloring Problem} (GCP) is the problem of attempting to color a set of bordering regions such that no region has the same color as its neighbors using three or four colors. For example consider the problem of coloring a map of the USA (Figure \ref{usa}), using only four colors and ensuring that no neighboring states share the same color. This is the motivation of the graph coloring problem. 
\begin{figure}[H]
	\centering
	\includegraphics[width=\linewidth]{../images/usa}
	\caption{Map of the United States of America satisfying the graph coloring problem.}
	\label{usa}
\end{figure}
It can be shown that the map coloring problem reduces to the graph coloring problem if we represent the states as the vertices of the graph, and the borders between states as the edges of the graph. This configuration produces a \textit{ planar} graph , a graph with no edge intersections. For this project, we will concentrate on solving the graph coloring problem for \textit{maximally planar} graphs (Figure \ref{simple}), planar graphs for which adding any addition edge would make the graph non-planar.
\begin{figure}[H]
	\centering
	\includegraphics[width=0.3\linewidth]{../images/simple_graph}
	\caption{Simple example of a maximally planar graph.}
	\label{simple}
\end{figure}
We are tasked with solving this problem five different ways: Minimum Conflicts, Simple Backtracking, Backtracking with Forward Checking, Backtracking with Constraint Propagation (MAC), and Local Search using a Genetic Algorithm. To test these algorithms, we will first need to build a problem generating program that can produce a random set of maximally planar graphs. Using the problem generator, we will calculate a set of graphs between the sizes \{10, 20, 30,...,100\} and then use the five graph coloring algorithms to solve the GCP. We will measure GCP algorithm performance by how many vertex colorings it requires to find a solution.  \par The solution to this project will be implemented in Python 3.4 using the graphics library to visualize the graphs.
\section{Problem Generation}
For this project, we will need to create maximally planar graphs (MPGs) from a set of randomly scattered points. For an arbitrary set of points, there is no unique MPG that can be constructed. To solve this issue, the problem statement has provided a prescription for calculating an MPG. 
\begin{quote}
Select some point $X$ at random and connect $X$ by a straight line
to the nearest point $Y$ such that $X$ is not already connected to $Y$ and line crosses no other line. Repeat the
previous step until no more connections are possible.
\end{quote}
To implement this algorithm we will first create a complete graph (graph where each point is connected to every other). Each vertex will have a list of associated edges, sorted by length. As instructed, we will then select a point at random and inspect the first unchecked edge $E$. If any edge crosses $E$ it is eliminated from the graph. We then repeat this process until every edge has been checked. \par
An algorithm to determine if two edges (or line segments) cross has been outlined by LaMothe \cite{tricks} and described in detail here. Let's start by defining two line segments 
\begin{align*}
	A = \{ \ve{a},\; \ve{a'} \} \\
	B = \{\ve{b}, \;\ve{b'}\}
\end{align*}
where $\ve{a}, \; \ve{a'}, \; \ve{b}, \text{ and } \ve{b'}$ are points on the ends of the line segments.

\section{Experiment Design}

\section{Solution Overview}

\section{GCP Algorithms}
\subsection{Minimum Conflicts}
Minimum conflicts is a local search algorithm that can be used to solve the graph coloring problem. The approach is to randomly color the graph and assign each node a conflict score determined by how many of its neighbors have a color in common with that node. The algorithm then randomly chooses a conflicting node and change it to a color that results in the least conflicts. The algorithm continues to choose conflicting nodes and changing them until there are no more conflicts or an until a predefined limit on how long the program runs is reached. \par
There are several design decisions that have been made to facilitate the use of this algorithm. One decision is to include a conflict score variable with each point. In addition, some functions will have to be implemented such as a node selector that ignores nodes with zero conflicts and a function that selects a color that causes the least amount of conflicts. The experimental design will consist of a counter that is incremented each time a new node is chosen. This should be sufficient as all other operations should be of constant or O(N) complexity at worst.
\subsection{Simple Backtracking}
Backtracking is a variant of the depth-first search. For the graph coloring problem the approach is to choose a node and assign it a valid color. The algorithm continues to do this until a solution is found or a node is reached that cannot be assigned a valid color. If the node cannot be assigned a valid color the algorithm back tracks to the previous node and tries a different color.  \par
This algorithm will have several design features that are required to make the algorithm work. The foremost feature is that the algorithm will be recursive as it is a convenient way of keeping track of old states. Each instance of the recursive function will have a list of previously used colors so it can know when to backtrack. The experimental design will consist of a counter that is incremented each time a node is colored. 
\subsection{Backtracking with Forward Checking}
\subsection{Backtracking with Constraint Propagation}
\subsection{Local Search using a Genetic Algorithm}


\end{multicols}

	%\bibliographystyle{apj}
	\bibliographystyle{unsrt}
	\bibliography{sources}
\end{document}